\subsection{根号を含む式への置換}

第3の方法では,関数の名前を使わずに
\[
t=x+\sqrt{1+x^2}
\]
と置く.この和を選ぶ理由は,差との積が$1$になるためである.
\[
\bigl(\sqrt{1+x^2}+x\bigr)
\bigl(\sqrt{1+x^2}-x\bigr)=1.
\]
したがって$\sqrt{1+x^2}-x=\frac1t$であり,この式と$t$の式の和と差を取って2で割れば
\[
x=\frac12\left(t-\frac1t\right),\qquad
\sqrt{1+x^2}=\frac12\left(t+\frac1t\right),\qquad
\dd x=\frac12\left(1+\frac1{t^2}\right)\dd t.
\]
$t>0$であり,この$x$の式は$t>0$から実数全体への狭義単調増加な対応である.根号も微分も有理式へ変わったので,
\[
\begin{aligned}
\int\sqrt{1+x^2}\dd x
&=\frac14\int\left(t+\frac2t+\frac1{t^3}\right)\dd t\\
&=\frac{t^2}{8}+\frac12\log t-\frac1{8t^2}+C.
\end{aligned}
\]
ここで,\,$x$と根号の$t$による式を掛けると
\[
x\sqrt{1+x^2}=\frac14\left(t^2-\frac1{t^2}\right)
\]
なので,\,$\frac{t^2}{8}-\frac1{8t^2}=\frac12x\sqrt{1+x^2}$である.また$\frac12\log t=\frac12\log\bigl(x+\sqrt{1+x^2}\bigr)$なので,式\eqref{式:一と二乗の根号積分}と同じ答えになる.

\paragraphbf{3つの置換のつながり}
3つの置換は,変数どうしの置換で直接結ばれている.まず双曲線置換と代数的な置換では,\,$x=\sinh u$なら
\[
t=x+\sqrt{1+x^2}=\sinh u+\cosh u=\e^u
\]
となる.つまり第3の方法は,双曲線の媒介変数$u$の代わりに$t=\e^u$を変数として,\,$\e^u$と$\e^{-u}$の式を有理式に直したものである.実際,\,$\cosh^2u=\frac14\left(t+\frac1t\right)^2$,\,$\dd u=\frac{\dd t}{t}$を代入すれば,第2の方法の$\int\cosh^2u\dd u$は第3の方法の$t$の積分になる.

次に三角置換と双曲線置換では,\,$x=\tan\theta=\sinh u$であり,
\[
\sec\theta=\sqrt{1+\tan^2\theta}=\sqrt{1+\sinh^2u}=\cosh u
\]
となる.\,$\dd x$を2通りに書くと$\sec^2\theta\dd\theta=\cosh u\dd u$なので
\[
\int\sec^3\theta\dd\theta
=\int\sec\theta\cdot\sec^2\theta\dd\theta
=\int\cosh u\cdot\cosh u\dd u
=\int\cosh^2u\dd u
\]
である.三角置換で現れた$\sec\theta$の3乗の積分は,双曲線置換の$\cosh u$の2乗の積分と同じものだったのだ.三角置換では$\dd x$が$\sec^2\theta\dd\theta$となって次数が1つ上がり,部分積分が必要になった.双曲線置換では$\dd x=\cosh u\dd u$の次数が根号と同じなので,2倍角の公式だけで済んだ.

さらに,三角置換の答えに現れた$\sec\theta+\tan\theta$は$\cosh u+\sinh u=\e^u=t$である.3つの変数は
\[
x=\tan\theta=\sinh u,\qquad
t=\sec\theta+\tan\theta=\e^u
\]
で互いに移り合う.\chapterref{章:部分積分}で,\,$\{x+\sqrt{1+x^2}\}'$の中に元の$x+\sqrt{1+x^2}$が現れたのも,これが$t=\e^u$であり,\,$\frac{\dd t}{\dd u}=t$となるからである.

最後に答えを微分して確認しておく.\,$s=\sqrt{1+x^2}$と書けば$s'=\frac xs$であり,
\[
\frac{\dd}{\dd x}\log(x+s)
=\cfrac{1+\cfrac xs}{x+s}=\frac1s.
\]
したがって式\eqref{式:一と二乗の根号積分}の右辺の微分は
\[
\frac12\left(s+\frac{x^2}{s}+\frac1s\right)
=\frac12(s+s)=\sqrt{1+x^2}
\]
となる.3つの計算は同じ原始関数へ到達した.



